% =====================================================================
%  Appendix F: Proof Sketches and Theorem Dependencies
% =====================================================================

\chapter{Proof Sketches and Theorem Dependencies}
\label{app:proof-sketches}

This appendix records why the main theorems are true and which earlier facts they use.
It is not a full proof chapter.  It is a map of the logic.

Each entry gives:

\[
\begin{array}{c|l}
\text{Item} & \text{Meaning} \\ \hline
\text{Statement} & \text{the theorem in compact form}\\[1mm]
\text{Hypotheses} & \text{the assumptions that make the theorem safe}\\[1mm]
\text{Proof idea} & \text{the main mechanism behind the proof}\\[1mm]
\text{Dependencies} & \text{earlier facts used}\\[1mm]
\text{What can fail} & \text{why the hypotheses are not decoration}
\end{array}
\]

\[
\begin{array}{c|c}
\text{Theorem} & \text{Main source of the proof} \\ \hline
\text{Differentiability from continuous partials} & \text{single-variable FTC along coordinate steps}\\[1mm]
\text{Equality of mixed partials} & \text{two orders of measuring rectangle change}\\[1mm]
\text{Fubini} & \text{slice sums}\\[1mm]
\text{Change of variables} & \text{local linearization and determinant scaling}\\[1mm]
\text{Green} & \text{small rectangles and cancellation}\\[1mm]
\text{Stokes} & \text{pull back to Green's theorem}\\[1mm]
\text{Divergence theorem} & \text{small boxes and cancellation}\\[1mm]
\text{Generalized Stokes} & \text{boundary cancellation in oriented cells}
\end{array}
\]

% ---------------------------------------------------------------------
\section{Differentiability and continuous partial derivatives}
\label{app:F.1}

\subsection*{Statement}

Let
\[
    f:\mathbb R^n\to\mathbb R.
\]
If the first partial derivatives of \(f\) exist in a neighborhood of \(\mathbf p\) and are
continuous at \(\mathbf p\), then \(f\) is differentiable at \(\mathbf p\).

In two variables:
\[
\boxed{
    f(a+h,b+k)
    =
    f(a,b)
    +
    f_x(a,b)h
    +
    f_y(a,b)k
    +
    \text{error}
}
\]
where
\[
\boxed{
    \frac{\text{error}}{\sqrt{h^2+k^2}}\to0
    \quad\text{as }(h,k)\to(0,0).
}
\]

\subsection*{Differential form}

\[
\boxed{
    df_{\mathbf p}
    =
    f_{x_1}(\mathbf p)\,dx_1+\cdots+f_{x_n}(\mathbf p)\,dx_n.
}
\]

In two variables:
\[
\boxed{
    df_{(a,b)}=f_x(a,b)\,dx+f_y(a,b)\,dy.
}
\]

\subsection*{Hypotheses}

\[
\begin{array}{ll}
1. & \text{The partial derivatives exist near the point.}\\[1mm]
2. & \text{The partial derivatives are continuous at the point.}
\end{array}
\]

The theorem is a sufficient condition, not a necessary condition.

A function can be differentiable even if its partial derivatives are not continuous.

\subsection*{Proof idea in two variables}

Move from
\[
    (a,b)
\]
to
\[
    (a+h,b+k)
\]
in two coordinate steps:
\[
    (a,b)\to(a+h,b)\to(a+h,b+k).
\]

Then
\[
\begin{aligned}
f(a+h,b+k)-f(a,b)
&=
\bigl[f(a+h,b+k)-f(a+h,b)\bigr]\\
&\quad+
\bigl[f(a+h,b)-f(a,b)\bigr].
\end{aligned}
\]

Apply the one-variable Mean Value Theorem or Fundamental Theorem of Calculus to each
coordinate slice:
\[
    f(a+h,b)-f(a,b)
    \approx
    f_x(a,b)h,
\]
\[
    f(a+h,b+k)-f(a+h,b)
    \approx
    f_y(a,b)k.
\]

Continuity of \(f_x\) and \(f_y\) makes the approximation uniform as
\[
    (h,k)\to(0,0).
\]

So the first-order part is
\[
\boxed{
    f_x(a,b)h+f_y(a,b)k.
}
\]

\subsection*{Dependencies}

\[
\begin{array}{ll}
1. & \text{Single-variable Mean Value Theorem or FTC.}\\[1mm]
2. & \text{Continuity of partial derivatives.}\\[1mm]
3. & \text{Linear approximation in one variable.}
\end{array}
\]

\subsection*{What can fail}

Partial derivatives can exist without differentiability.

Define
\[
    f(x,y)=
    \begin{cases}
    \dfrac{xy}{\sqrt{x^2+y^2}}, & (x,y)\ne(0,0),\\[2mm]
    0, & (x,y)=(0,0).
    \end{cases}
\]

At the origin,
\[
    f_x(0,0)=0,
    \qquad
    f_y(0,0)=0.
\]

If these partial derivatives gave a true differential, the linear approximation would be
zero.

But along the path
\[
    (x,y)=(t,t),
\]
we get
\[
    f(t,t)=\frac{t^2}{\sqrt{2t^2}}=\frac{|t|}{\sqrt2}.
\]

The distance to the origin is
\[
    \sqrt{t^2+t^2}=\sqrt2\,|t|.
\]

The ratio is
\[
    \frac{|t|/\sqrt2}{\sqrt2\,|t|}=\frac12,
\]
which does not go to \(0\).

So \(f\) is not differentiable at the origin.

\subsection*{Moral}

Coordinate slopes alone do not guarantee a true tangent plane.  Continuous partial
derivatives stitch the coordinate information into one linear approximation.

% ---------------------------------------------------------------------
\section{Equality of mixed partials}
\label{app:F.2}

\subsection*{Statement}

If \(f_{xy}\) and \(f_{yx}\) are continuous near a point, then
\[
\boxed{
    f_{xy}=f_{yx}
}
\]
at that point.

More generally, if the relevant higher partial derivatives are continuous, then the order
of differentiation can be switched.

\subsection*{Hypotheses}

\[
\begin{array}{ll}
1. & f_{xy}\text{ and }f_{yx}\text{ exist near the point.}\\[1mm]
2. & f_{xy}\text{ and }f_{yx}\text{ are continuous near the point.}
\end{array}
\]

\subsection*{Proof idea}

Measure the net change of \(f\) around a small rectangle:
\[
    [a,a+h]\times[b,b+k].
\]

Form the corner difference
\[
\boxed{
    \Delta
    =
    f(a+h,b+k)-f(a+h,b)-f(a,b+k)+f(a,b).
}
\]

Compute \(\Delta\) in two ways.

First, take \(x\)-changes, then \(y\)-changes.  This connects \(\Delta\) to \(f_{xy}\).

Second, take \(y\)-changes, then \(x\)-changes.  This connects \(\Delta\) to \(f_{yx}\).

Both computations measure the same corner difference.  After division by
\[
    hk
\]
and letting
\[
    h,k\to0,
\]
continuity forces
\[
    f_{xy}(a,b)=f_{yx}(a,b).
\]

\subsection*{Dependencies}

\[
\begin{array}{ll}
1. & \text{Single-variable FTC or Mean Value Theorem.}\\[1mm]
2. & \text{Continuity of the mixed partial derivatives.}\\[1mm]
3. & \text{Rectangle difference quotients.}
\end{array}
\]

\subsection*{What can fail}

If the mixed partials are not continuous, they can differ.

Define
\[
    f(x,y)=
    \begin{cases}
    \dfrac{xy(x^2-y^2)}{x^2+y^2}, & (x,y)\ne(0,0),\\[2mm]
    0, & (x,y)=(0,0).
    \end{cases}
\]

Then
\[
    f_x(0,y)=-y,
\]
so
\[
    f_{xy}(0,0)=-1.
\]

Also,
\[
    f_y(x,0)=x,
\]
so
\[
    f_{yx}(0,0)=1.
\]

Thus
\[
\boxed{
    f_{xy}(0,0)\ne f_{yx}(0,0).
}
\]

\subsection*{Moral}

The equality of mixed partials is not only about the symbols \(xy\) and \(yx\).  It is
about enough smoothness to make the small-rectangle measurement independent of the
order in which it is taken.

% ---------------------------------------------------------------------
\section{Fubini's theorem and change of variables}
\label{app:F.3}

\subsection*{Fubini's theorem: statement}

If \(f\) is continuous on a rectangle
\[
    R=[a,b]\times[c,d],
\]
then
\[
\boxed{
    \iint_R f(x,y)\,dA
    =
    \int_a^b\int_c^d f(x,y)\,dy\,dx
    =
    \int_c^d\int_a^b f(x,y)\,dx\,dy.
}
\]

For a solid box,
\[
\boxed{
    \iiint_B f\,dV
}
\]
can be computed by any order of iterated integration, as long as the corresponding
bounds describe the same box.

\subsection*{Fubini: hypotheses}

For the version used most often in this book:

\[
\begin{array}{ll}
1. & \text{The region is a rectangle, box, or a region described by clean bounds.}\\[1mm]
2. & \text{The function is continuous, or at least integrable in a safe sense.}
\end{array}
\]

\subsection*{Fubini: proof idea}

A double integral is a limit of sums over small rectangles:
\[
    \sum f(x_i^*,y_j^*)\,\Delta x\,\Delta y.
\]

Group the same sum by columns:
\[
    \sum_i
    \left(
        \sum_j f(x_i^*,y_j^*)\,\Delta y
    \right)
    \Delta x.
\]

The inner sum becomes an integral in \(y\).  The outer sum becomes an integral in
\(x\).

Group by rows instead, and the order reverses.

\subsection*{Fubini: dependencies}

\[
\begin{array}{ll}
1. & \text{Riemann sums.}\\[1mm]
2. & \text{Single-variable integration.}\\[1mm]
3. & \text{Continuity or integrability.}
\end{array}
\]

\subsection*{Fubini: what can fail}

For ordinary continuous functions on bounded regions, Fubini is safe.

For improper integrals, changing the order can change the value if the integral is not
absolutely convergent.

The safe rule is:

\[
\boxed{
    \text{If }\iint |f|\,dA\text{ is finite, changing order is safe.}
}
\]

Without absolute convergence, be careful.

\subsection*{Change of variables: statement}

For a \(C^1\) coordinate map
\[
    T(u,v)=(x(u,v),y(u,v)),
\]
the change of variables formula is
\[
\boxed{
    \iint_R f(x,y)\,dA
    =
    \iint_{\widetilde R}
    f(T(u,v))
    \left|
        \det DT(u,v)
    \right|
    \,du\,dv.
}
\]

In differential forms language:
\[
\boxed{
    T^*(dx\wedge dy)
    =
    \det DT\,du\wedge dv.
}
\]

So the oriented form is
\[
\boxed{
    \iint_R f\,dx\wedge dy
    =
    \iint_{\widetilde R}
    f(T(u,v))\det DT\,du\wedge dv.
}
\]

The ordinary area formula uses
\[
    |\det DT|.
\]

\subsection*{Change of variables: hypotheses}

The standard version assumes:

\[
\begin{array}{ll}
1. & T\text{ is }C^1.\\[1mm]
2. & T\text{ is one-to-one on the region, except possibly on boundary pieces.}\\[1mm]
3. & \det DT\ne0\text{ in the interior, except possibly on harmless boundary pieces.}\\[1mm]
4. & \text{The regions are reasonably nice.}
\end{array}
\]

\subsection*{Change of variables: proof idea}

Near a point \((u,v)\), the map \(T\) is almost linear:
\[
    T(u+\Delta u,v+\Delta v)
    \approx
    T(u,v)+DT
    \begin{pmatrix}
    \Delta u\\
    \Delta v
    \end{pmatrix}.
\]

A tiny rectangle in the \(uv\)-plane becomes almost a parallelogram in the \(xy\)-plane.

The area scale factor of that parallelogram is
\[
\boxed{
    |\det DT|.
}
\]

Add over many tiny rectangles and pass to the limit.

\subsection*{Change of variables: dependencies}

\[
\begin{array}{ll}
1. & \text{Differentiability as local linear approximation.}\\[1mm]
2. & \text{Determinant as area or volume scale factor.}\\[1mm]
3. & \text{Riemann sums.}\\[1mm]
4. & \text{One-to-one coordinate descriptions.}
\end{array}
\]

\subsection*{Change of variables: what can fail}

If the map is not one-to-one, the integral may count points more than once.

Example: polar coordinates with
\[
    0\le \theta\le4\pi
\]
trace a full disk twice.

If the Jacobian vanishes on a whole region, the map may collapse area.

If the orientation matters, do not replace \(\det DT\) by \(|\det DT|\).  The sign is part
of the form.

\subsection*{Moral}

Fubini explains how to add by slices.  Change of variables explains how small pieces
change size and orientation under a new coordinate language.

% ---------------------------------------------------------------------
\section{Green's theorem}
\label{app:F.4}

\subsection*{Statement in forms}

Let
\[
    \omega=P\,dx+Q\,dy.
\]

If \(R\) is a positively oriented plane region, then
\[
\boxed{
    \int_{\partial R}\omega
    =
    \int_R d\omega.
}
\]

Since
\[
    d\omega=(Q_x-P_y)\,dx\wedge dy,
\]
this is
\[
\boxed{
    \int_{\partial R}P\,dx+Q\,dy
    =
    \iint_R (Q_x-P_y)\,dx\wedge dy.
}
\]

\subsection*{Classical statement}

For
\[
    \mathbf F=(P,Q),
\]
\[
\boxed{
    \oint_{\partial R}\mathbf F\cdot d\mathbf r
    =
    \iint_R
    \left(
        \frac{\partial Q}{\partial x}
        -
        \frac{\partial P}{\partial y}
    \right)
    dA.
}
\]

\subsection*{Hypotheses}

\[
\begin{array}{ll}
1. & R\text{ is a bounded plane region with piecewise smooth boundary.}\\[1mm]
2. & \partial R\text{ has positive orientation.}\\[1mm]
3. & P\text{ and }Q\text{ have continuous first partial derivatives on a region containing }R.
\end{array}
\]

Positive orientation means:

\[
\boxed{
    \text{walk so the region stays on your left.}
}
\]

\subsection*{Proof idea}

First prove the theorem for one small rectangle.

For
\[
    R=[a,a+h]\times[b,b+k],
\]
the line integral around the rectangle compares opposite sides.

Horizontal sides produce the term
\[
    -P_y\,hk.
\]

Vertical sides produce the term
\[
    Q_x\,hk.
\]

Together:
\[
    \int_{\partial R}P\,dx+Q\,dy
    \approx
    (Q_x-P_y)hk.
\]

Then subdivide a larger region into many small rectangles.

Interior edges are walked twice in opposite directions, so their contributions cancel.

Only the outer boundary remains.

Let the mesh size go to zero.

\subsection*{Dependencies}

\[
\begin{array}{ll}
1. & \text{Line integrals of \(1\)-forms.}\\[1mm]
2. & \text{Exterior derivative }d(P\,dx+Q\,dy).\\[1mm]
3. & \text{Fubini's theorem.}\\[1mm]
4. & \text{Cancellation of oppositely oriented interior edges.}
\end{array}
\]

\subsection*{What can fail}

Orientation matters.

If the boundary is walked clockwise instead of counterclockwise, the sign changes.

The field must be defined on the whole region.

Example:
\[
    \omega=
    \frac{-y}{x^2+y^2}\,dx
    +
    \frac{x}{x^2+y^2}\,dy.
\]

This form is defined on
\[
    \mathbb R^2\setminus\{(0,0)\}.
\]

On that punctured plane,
\[
    d\omega=0.
\]

But around the unit circle,
\[
    \int_{C}\omega=2\pi.
\]

Green's theorem cannot be applied to the disk bounded by \(C\), because \(\omega\) is
not defined at the origin.

\subsection*{Moral}

Green's theorem is a cancellation theorem.  Local rotation inside a region adds up to
circulation around the boundary, as long as the field is smooth on the whole region.

% ---------------------------------------------------------------------
\section{Stokes' theorem}
\label{app:F.5}

\subsection*{Statement in forms}

Let
\[
    \omega=P\,dx+Q\,dy+R\,dz.
\]

If \(S\) is an oriented surface with oriented boundary \(\partial S\), then
\[
\boxed{
    \int_{\partial S}\omega
    =
    \int_S d\omega.
}
\]

\subsection*{Classical statement}

For
\[
    \mathbf F=(P,Q,R),
\]
\[
\boxed{
    \oint_{\partial S}\mathbf F\cdot d\mathbf r
    =
    \iint_S
    (\nabla\times\mathbf F)\cdot\widehat{\mathbf n}\,dS.
}
\]

\subsection*{Hypotheses}

\[
\begin{array}{ll}
1. & S\text{ is an oriented piecewise smooth surface.}\\[1mm]
2. & \partial S\text{ has the boundary orientation induced by }S.\\[1mm]
3. & P,Q,R\text{ have continuous first partial derivatives near }S.\\[1mm]
4. & \omega\text{ is defined on a region containing }S.
\end{array}
\]

Boundary orientation is given by the right-hand rule:

\[
\boxed{
    \text{thumb points along the chosen normal, fingers curl along }\partial S.
}
\]

\subsection*{Proof idea}

Parametrize the surface:
\[
    \mathbf r(u,v).
\]

Pull back the \(1\)-form:
\[
    \mathbf r^*\omega.
\]

Now \(\mathbf r^*\omega\) is a \(1\)-form in the \(uv\)-plane.

Apply Green's theorem in the parameter region:
\[
    \int_{\partial D}\mathbf r^*\omega
    =
    \int_D d(\mathbf r^*\omega).
\]

Use the pullback identity:
\[
\boxed{
    d(\mathbf r^*\omega)=\mathbf r^*(d\omega).
}
\]

This gives
\[
    \int_{\partial S}\omega
    =
    \int_S d\omega.
\]

\subsection*{Dependencies}

\[
\begin{array}{ll}
1. & \text{Green's theorem.}\\[1mm]
2. & \text{Parametrized surfaces.}\\[1mm]
3. & \text{Pullbacks of forms.}\\[1mm]
4. & \text{Exterior derivative commutes with pullback.}\\[1mm]
5. & \text{Orientation from }\mathbf r_u\times\mathbf r_v.
\end{array}
\]

\subsection*{What can fail}

Orientation matters.

If the surface normal is reversed, then the boundary orientation reverses and both
sides change sign.

The field must be defined on the surface.

Singularities can break the theorem if they lie on the chosen surface.

If \(S\) has no boundary, then
\[
    \partial S=\varnothing,
\]
so Stokes gives
\[
\boxed{
    \int_S d\omega=0.
}
\]

Classically:
\[
\boxed{
    \iint_S(\nabla\times\mathbf F)\cdot\widehat{\mathbf n}\,dS=0
}
\]
for a closed oriented surface, provided \(\mathbf F\) is smooth near \(S\).

\subsection*{Moral}

Stokes' theorem is Green's theorem on a curved surface.  The pullback flattens the
surface into a parameter region where Green's theorem can do the work.

% ---------------------------------------------------------------------
\section{Divergence theorem}
\label{app:F.6}

\subsection*{Statement in forms}

Let
\[
    \Phi_{\mathbf F}
    =
    P\,dy\wedge dz
    +
    Q\,dz\wedge dx
    +
    R\,dx\wedge dy.
\]

If \(E\) is an oriented solid with outward boundary \(\partial E\), then
\[
\boxed{
    \int_{\partial E}\Phi_{\mathbf F}
    =
    \int_E d\Phi_{\mathbf F}.
}
\]

Since
\[
    d\Phi_{\mathbf F}
    =
    (P_x+Q_y+R_z)\,dx\wedge dy\wedge dz,
\]
this becomes
\[
\boxed{
    \int_{\partial E}\Phi_{\mathbf F}
    =
    \iiint_E
    (P_x+Q_y+R_z)\,dx\wedge dy\wedge dz.
}
\]

\subsection*{Classical statement}

\[
\boxed{
    \iint_{\partial E}
    \mathbf F\cdot\widehat{\mathbf n}\,dS
    =
    \iiint_E
    \nabla\cdot\mathbf F\,dV.
}
\]

\subsection*{Hypotheses}

\[
\begin{array}{ll}
1. & E\text{ is a bounded solid with piecewise smooth boundary.}\\[1mm]
2. & \partial E\text{ is oriented outward.}\\[1mm]
3. & \mathbf F\text{ has continuous first partial derivatives on a region containing }E.
\end{array}
\]

\subsection*{Proof idea}

First prove the theorem for a small rectangular box.

Flux through the two \(x\)-faces gives approximately
\[
    P_x\,\Delta x\,\Delta y\,\Delta z.
\]

Flux through the two \(y\)-faces gives approximately
\[
    Q_y\,\Delta x\,\Delta y\,\Delta z.
\]

Flux through the two \(z\)-faces gives approximately
\[
    R_z\,\Delta x\,\Delta y\,\Delta z.
\]

Together:
\[
    \text{outward flux}
    \approx
    (P_x+Q_y+R_z)\,\Delta V.
\]

Now subdivide a solid into many small boxes.

Flux through interior faces cancels because each shared face is outward for one box and
inward for the neighboring box.

Only the outer boundary remains.

Let the box sizes shrink.

\subsection*{Dependencies}

\[
\begin{array}{ll}
1. & \text{Flux forms.}\\[1mm]
2. & \text{Exterior derivative of a \(2\)-form.}\\[1mm]
3. & \text{Triple integrals.}\\[1mm]
4. & \text{Cancellation across shared faces.}\\[1mm]
5. & \text{Single-variable FTC along coordinate directions.}
\end{array}
\]

\subsection*{What can fail}

The field must be defined inside the solid, not only on the boundary.

Example:
\[
    \mathbf F(x,y,z)
    =
    \frac{(x,y,z)}{(x^2+y^2+z^2)^{3/2}}.
\]

This field is not defined at the origin.

Away from the origin,
\[
    \nabla\cdot\mathbf F=0.
\]

But on the unit sphere,
\[
    \mathbf F=\widehat{\mathbf n},
\]
so
\[
    \iint_S \mathbf F\cdot\widehat{\mathbf n}\,dS
    =
    \iint_S 1\,dS
    =
    4\pi.
\]

The divergence theorem cannot be applied to the unit ball because the field is singular
at the origin.

\subsection*{Moral}

The divergence theorem is a flux bookkeeping theorem.  What leaves through the
boundary is accounted for by sources and sinks inside, provided the field is smooth
throughout the solid.

% ---------------------------------------------------------------------
\section{Generalized Stokes theorem}
\label{app:F.7}

\subsection*{Statement}

For an oriented \(k\)-dimensional region \(M\) and a \((k-1)\)-form \(\omega\),
\[
\boxed{
    \int_{\partial M}\omega
    =
    \int_M d\omega.
}
\]

This single statement includes:

\[
\begin{array}{c|c|c}
\text{Theorem} & M & \omega \\ \hline
\text{FTC} & \text{interval} & 0\text{-form}\\[1mm]
\text{Green} & \text{plane region} & 1\text{-form}\\[1mm]
\text{Stokes} & \text{surface} & 1\text{-form}\\[1mm]
\text{Divergence} & \text{solid} & 2\text{-form}
\end{array}
\]

\subsection*{Degree matching}

If
\[
    \dim M=k,
\]
then
\[
    \omega
\]
must be a \((k-1)\)-form and
\[
    d\omega
\]
is a \(k\)-form.

\[
\boxed{
    \int_{\partial M}\omega
}
\]
makes sense because
\[
    \dim(\partial M)=k-1.
\]

\[
\boxed{
    \int_M d\omega
}
\]
makes sense because \(d\omega\) has degree \(k\).

\subsection*{Hypotheses}

The clean version assumes:

\[
\begin{array}{ll}
1. & M\text{ is oriented and piecewise smooth.}\\[1mm]
2. & \partial M\text{ has the induced boundary orientation.}\\[1mm]
3. & \omega\text{ is smooth enough on a region containing }M.\\[1mm]
4. & \text{The integrals are defined.}
\end{array}
\]

\subsection*{Proof idea}

Break \(M\) into many small oriented pieces.

On each small piece, the theorem reduces to the coordinate version.

When the small pieces are added, every interior boundary appears twice with opposite
orientations.

Those interior contributions cancel.

Only the boundary of the whole region remains.

This is the same cancellation pattern behind Green's theorem and the divergence theorem.

\subsection*{Dependencies}

\[
\begin{array}{ll}
1. & \text{The one-dimensional FTC on coordinate intervals.}\\[1mm]
2. & \text{Wedge products and orientation.}\\[1mm]
3. & \text{Exterior derivative.}\\[1mm]
4. & \text{Pullbacks to parametrized pieces.}\\[1mm]
5. & \text{Boundary orientation conventions.}\\[1mm]
6. & \text{Cancellation of interior boundaries.}
\end{array}
\]

\subsection*{Boundary of a boundary}

A boundary has no boundary:
\[
\boxed{
    \partial(\partial M)=0.
}
\]

The form version is
\[
\boxed{
    d^2=0.
}
\]

These two facts match.

If
\[
    \partial(\partial M)=0,
\]
then applying Stokes twice gives
\[
    \int_{\partial(\partial M)}\omega=0.
\]

On the derivative side,
\[
    \int_M d(d\omega)=\int_M 0=0.
\]

\subsection*{Closed and exact}

A form is closed if
\[
\boxed{
    d\omega=0.
}
\]

A form is exact if
\[
\boxed{
    \omega=d\eta.
}
\]

Exact implies closed:
\[
\boxed{
    \omega=d\eta
    \quad\Longrightarrow\quad
    d\omega=d^2\eta=0.
}
\]

Closed need not imply exact.

A hole in the domain can prevent exactness.

\subsection*{What can fail}

If the form is not defined on all of \(M\), Stokes' theorem may fail.

If the orientation of \(\partial M\) is wrong, the sign is wrong.

If \(M\) has corners, edges, or several patches, it must be handled as a piecewise
smooth object with consistent orientations.

If singularities are inside the region, remove small neighborhoods around them and
include the new boundary pieces.

\subsection*{Moral}

Every version of Stokes' theorem says the same thing:

\[
\boxed{
    \text{local derivative information inside}
    =
    \text{accumulated boundary information}.
}
\]

The shape of the region changes from interval to surface to solid, but the mechanism is
always the same: orient, differentiate, integrate, and cancel interior boundaries.

% ---------------------------------------------------------------------
\section*{Dependency chart}

\[
\begin{array}{c|l}
\text{Result} & \text{Depends on} \\ \hline
\text{Differentiability from continuous partials} &
\text{one-variable FTC/MVT, continuity}\\[1mm]

\text{Equality of mixed partials} &
\text{rectangle differences, FTC/MVT, continuity}\\[1mm]

\text{Fubini} &
\text{Riemann sums, single-variable integration}\\[1mm]

\text{Change of variables} &
\text{differentiability, determinants, Riemann sums}\\[1mm]

\text{Green} &
\text{small rectangles, Fubini, cancellation}\\[1mm]

\text{Stokes} &
\text{Green, parametrized surfaces, pullbacks}\\[1mm]

\text{Divergence theorem} &
\text{small boxes, flux, cancellation}\\[1mm]

\text{Generalized Stokes} &
\text{orientation, exterior derivative, pullbacks, cancellation}
\end{array}
\]

% ---------------------------------------------------------------------
\section*{One-page theorem summary}

\[
\boxed{
    \text{continuous partials near }\mathbf p
    \Longrightarrow
    \text{differentiability at }\mathbf p
}
\]

\[
\boxed{
    f_{xy}=f_{yx}
    \quad\text{if the mixed partials are continuous}
}
\]

\[
\boxed{
    \iint_R f\,dA
    =
    \int\!\!\int f\,dy\,dx
    =
    \int\!\!\int f\,dx\,dy
}
\]
for continuous \(f\) on a suitable region.

\[
\boxed{
    dA=
    \left|
    \det DT
    \right|
    du\,dv
}
\]

\[
\boxed{
    dx\wedge dy
    =
    \det DT\,du\wedge dv
}
\]

\[
\boxed{
    \int_{\partial R}P\,dx+Q\,dy
    =
    \iint_R(Q_x-P_y)\,dx\wedge dy
}
\]

\[
\boxed{
    \oint_{\partial S}\mathbf F\cdot d\mathbf r
    =
    \iint_S(\nabla\times\mathbf F)\cdot\widehat{\mathbf n}\,dS
}
\]

\[
\boxed{
    \iint_{\partial E}\mathbf F\cdot\widehat{\mathbf n}\,dS
    =
    \iiint_E\nabla\cdot\mathbf F\,dV
}
\]

\[
\boxed{
    \int_{\partial M}\omega=\int_M d\omega
}
\]

\[
\boxed{
    d^2=0
}
\]

\[
\boxed{
    \partial(\partial M)=0
}
\]